% \section{Discussion and Future Work}
% \label{sec:future}

% \sysname{} has proven effective in improving the detection of routing device misconfigurations, but there are several avenues for further exploration and refinement. One challenge is expanding the framework to better handle large-scale, distributed networks with interdependent routers. These environments often involve complex relationships between devices, and future work could focus on extending \sysname{} to capture and analyze multi-device contexts by default. Incorporating cross-router consistency checks and advanced network-wide policy validation could be valuable enhancements.

% Another area for improvement is reducing false positives while maintaining high detection accuracy. Without a definitive ground truth, distinguishing between true misconfigurations and intentional deviations is inherently challenging. Future work could explore techniques such as refining prompt engineering, incorporating confidence scoring, or leveraging ensemble methods to improve robustness and reduce unnecessary alerts.

% Additionally, enhancing the LLM's ability to process real-time configuration changes is crucial. As networks evolve dynamically, incorporating real-time monitoring data into the context mining process would allow for continuous verification and faster misconfiguration detection. Integrating \sysname{} with network management tools that track operational metrics could provide valuable runtime data, further improving detection accuracy by correlating live data with configuration snapshots.
% Future enhancements can also involve proposing and implementing automated fixes for detected misconfigurations with real-time updates. This would help network operators quickly address issues, streamlining troubleshooting and repairs.


\section{Discussion and Future Work}
\label{sec:future}
A key insight from our work is that the choice of the underlying LLM is not trivial; it reveals a fundamental trade-off between precision and recall. This finding has significant implications for network operators, suggesting that the "best" model may depend on the task—a deep, exploratory audit might favor a high-recall model, while continuous, automated monitoring would benefit from a high-precision one.
A second key insight is that our interactive "ask-for-context" framework is more than a verifier; it is also a diagnostic tool. The context a model requests reveals its underlying reasoning strategy---whether it focuses on local semantics or global patterns. This makes our interactive approach not only more efficient than force-feeding all available context, but also more transparent, providing a crucial window into the model's "thought process" that is essential for building trust in its conclusions.

Building on these insights, future work can proceed in several directions. First, given the precision-recall trade-off, an exciting avenue is to develop an ensemble or meta-LLM framework that could dynamically choose the best model (or combination of models) for a given configuration line or verification task. Second, while our context mining is comprehensive, it could be enriched by incorporating live network state data from sources like BGP routing tables or SNMP, allowing for verification not just of the configuration's snapshots, but also its alignment with the network's real-time operational reality. Finally, the structured output of \sysname{} provides a strong foundation for automated remediation, where the framework could not only detect a misconfiguration but also propose and, with operator approval, apply a validated fix.